\documentclass[11pt]{article}
\usepackage[margin=1in]{geometry}
\usepackage{amsmath,amssymb,amsthm,booktabs,array}
\usepackage[T1]{fontenc}
\usepackage{lmodern}
\usepackage[hidelinks]{hyperref}
\newtheorem{theorem}{Theorem}
\newtheorem{corollary}[theorem]{Corollary}
\newcommand{\Sub}{\operatorname{Sub}}
\newcommand{\Hull}{\operatorname{Hull}}
\newcommand{\gen}{\operatorname{gen}}
\newcommand{\id}{\operatorname{id}}
\title{Correcting Klanderman et al.'s isolated-top transfer-system conjecture\\
via DeMark et al.'s tulip theorem}
\author{Haobo Ma\\\small ChronoAI Pte Ltd, Singapore
\and Wenlin Zhang\\\small National University of Singapore\\\small ChronoAI Pte Ltd, Singapore}
\date{November 14, 2026}
\begin{document}
\maketitle
\begin{abstract}
Conjecture 4.2 of Klanderman, Lewis, Monson, Shibata and Van Niel, in
\emph{Characterizing Transfer Systems for Non-Abelian Groups},
arXiv:2511.13439v2, Section 4, printed page 15, receives an explicit
correction by an elementary derivation from DeMark et al.'s tulip theorem
(arXiv:2604.00335v1, 1 April 2026, Proposition 3.17) and the hull theorem
of Mazur et al.\ (2024, Proposition 2.14).
Under the conjecture's two-component hypothesis with isolated top,
the transfer system is non-LSP precisely when $|G|$ is composite, and its
ordinary generated extension is compatible precisely when the original
system is saturated.
The non-LSP classification needs only top isolation, not two components;
under the original hypotheses, all compatible partners form the lattice of
conjugation-invariant order ideals of proper subgroups.
On $Q_8$, the system consisting of identities and
$e\to H$ for every proper subgroup $H$ violates compatibility: the generated
extension omits $\{\pm1\}\to\langle i\rangle$.
The classification also identifies the least compatible extension containing
$e\to G$ and separates ordinary transfer generation from compatibility closure.
\end{abstract}

\section{The statement and conventions}
Throughout, $G$ is a nontrivial finite group, $e$ denotes its trivial subgroup,
and relations are on all actual subgroups, not conjugacy classes.
A $G$-transfer system $T$ is a partial order refining subgroup inclusion,
containing identities, closed under composition, restriction
\[
 K\to H,\quad L\leq H\quad\Longrightarrow\quad K\cap L\to L,
\]
and simultaneous $G$-conjugation of source and target
\cite[Definition 2.1]{K}.
Connected components refer to the undirected graph of nonidentity arrows.

We use the ordered compatibility convention of
\cite[Definition 2.14]{K}: $(T,U)$ is compatible when $T\subseteq U$ and
\begin{equation}\label{eq:compatibility}
 B,C\leq A,\quad (B\to A)\in T,\quad (B\cap C\to B)\in U
 \quad\Longrightarrow\quad (C\to A)\in U.
\end{equation}
A transfer system is saturated if $L\leq K\leq H$ and $L\to H$ imply
$K\to H$. Its saturated hull is the least saturated extension.
It is \emph{lesser simply paired} (LSP) if every compatible greater system
is either its saturated hull or the complete inclusion relation
\cite[Definitions 2.15--2.17]{K}.
A system is \emph{cosaturated} if it is generated by arrows with target $G$,
and \emph{bisaturated} if it is both saturated and cosaturated
\cite[Definition 2.18 and Remark 2.19]{K}.
These are the notions studied by Mazur et al.\ \cite{M}.
Their Proposition 2.14 describes the saturated hull through connected
components, and their Lemma 4.14 already gives non-LSP for the isolated-top
two-component case on $C_{p^r q^s}$, for distinct primes and $r,s\geq1$.
DeMark et al.'s tulip theorem \cite[Proposition 3.17]{D} provides the
arbitrary-group compatibility statement underlying the correction below.
Our purpose is to give its explicit consequences for Conjecture 4.2,
including ordinary generation and the removal of the component-count
condition from the non-LSP conclusion, with self-contained proofs.
For cyclic prime-power groups, compatibility is already characterized by
\cite[Proposition 3.2]{HML}, and the LSP classification is already in
\cite[Remark 4.18]{M}.

The hypotheses of \cite[Conjecture 4.2, p.~15]{K} prescribe precisely
the two graph components represented by $e$ and $G$, with the $G$-component
consisting of the isolated vertex and its identity arrow.
Its asserted conclusion is:
\begin{quote}\small
then $T$ is not LSP and is compatible with
the transfer system generated by adding the edge $e\to G$ to $T$.
\end{quote}
Thus its conclusion is the conjunction
\begin{equation}\label{eq:conjecture}
 T\text{ is not LSP}\qquad\text{and}\qquad
 (T,\gen(T\cup\{e\to G\}))\text{ is compatible}.
\end{equation}
Here $\gen$ means the least ordinary transfer system containing its
arguments. In particular, it does not include a compatibility-closure step.
Mazur et al.\ explicitly distinguish that additional step in
\cite[Lemma 3.6]{M}.

\section{A classification under the component hypotheses}
Write $\mathcal C=\{(K,H):K\leq H\leq G\}$ for the complete system and
$\id=\{(H,H):H\leq G\}$ for the identity relation, and set
\begin{align*}
 P&=\id\cup\{(e,H):H<G\},\\
 V&=\id\cup\{(K,H):K\leq H<G\},\\
 W&=V\cup\{(e,G)\}.
\end{align*}
Thus $V$ is complete below $G$, with $G$ isolated. All three are transfer
systems. For $P$, a nonidentity arrow has trivial source and proper target;
restriction preserves both conditions, conjugation preserves properness,
and a composite either removes an identity or retains trivial source and
proper target. For $V$, an arrow with proper target stays so under
restriction and conjugation; a composite has proper final target unless
both arrows are the identity at $G$. For $W$, restriction of the new arrow
$e\to G$ gives $e\to H$ already in $V$. Composing this arrow can only
compose an identity at one end, and conjugation fixes it.

The system $V$ is exactly the \emph{tulip} $\mathrm{Tu}(G)$ of
\cite[Definition 3.15]{D}. In their additive-first pair convention,
\cite[Propositions 2.21 and 3.17]{D} say that an incomplete system
$U\supseteq V$ has maximal compatible lesser system $V$.
Thus every $T\subseteq V$ is compatible with such a $U$; the complete
$U$ is always compatible by \cite[Proposition 2.17]{D}.
Conversely, the hull results \cite[Propositions 2.14 and 3.5]{M} force
$V\subseteq U$ under the component hypotheses. The compatibility criterion
in the next theorem is therefore a specialization of these earlier results.
The proof also derives it directly from the axioms.

\begin{theorem}\label{thm:classification}
Suppose $T$ satisfies the component hypotheses of Conjecture 4.2.
Then $P\subseteq T\subseteq V$, and the following statements hold.
\begin{enumerate}
\item $\gen(T\cup\{e\to G\})=T\cup\{(e,G)\}$.
\item $\Hull(T)=V$. For a transfer system $U$, $(T,U)$ is compatible
if and only if $V\subseteq U$.
\item The ordinary generated extension in (1) is compatible with $T$
if and only if $T=V$, equivalently if and only if $T$ is saturated.
The least compatible extension containing $e\to G$ is $W$.
\item $T$ is LSP if and only if $|G|$ is prime.
\end{enumerate}
Conversely, every transfer system between $P$ and $V$ satisfies the
component hypotheses.
\end{theorem}

\begin{proof}
First, for an arrow $K\to H$ in any transfer system,
$e\to K$ implies $e\to H$ by composition, whereas $e\to H$
implies $e\to K$ by restriction along $K\leq H$.
It follows along undirected paths that the component containing $e$
consists exactly of the subgroups $H$ for which $e\to H$.
The hypotheses place every proper subgroup in this component, so
$P\subseteq T$. Isolation of $G$ gives $T\subseteq V$.
Conversely, these inclusions join every proper subgroup to $e$ and leave
$G$ isolated, giving exactly the required two components.

Let $U_0=T\cup\{(e,G)\}$. The only new arrow is fixed by conjugation.
Its restrictions $e\to H$ already belong to $T$ for $H<G$.
An arrow composing before it must end at $e$ and is an identity.
An arrow composing after it must start at $G$ and is an identity.
Thus $U_0$ is a transfer system. It contains the generators and is
contained in every transfer system containing them, proving (1).

Saturation of $e\to H$ for each $H<G$ forces every $K\to H$ with
$K\leq H$. Hence every saturated extension of $T$ contains $V$.
The system $V$ itself is saturated: an arrow with proper target forces
only arrows with that same proper target, and the identity at $G$
forces no new arrow. Therefore $\Hull(T)=V$.

If $(T,U)$ is compatible, take $B=e$ and any $C\leq A<G$ in
\eqref{eq:compatibility}. The arrow $e\to A$ lies in $T$ and
$B\cap C\to B$ is $e\to e$, so $C\to A$ must lie in $U$.
Consequently $V\subseteq U$.
Conversely, suppose $V\subseteq U$. Then $T\subseteq U$.
For $A<G$ the conclusion of \eqref{eq:compatibility} belongs to $V$.
For $A=G$, the antecedent $B\to G$ in $T$ forces $B=G$;
the other antecedent is then exactly the desired arrow $C\to G$ in $U$.
This proves (2).

By (1)--(2), compatibility of $U_0$ is equivalent to
$V\subseteq T\cup\{(e,G)\}$. All arrows of $V$ other than the
identity at $G$ have proper target, so this is equivalent to $T=V$.
Since $\Hull(T)=V$, it is also equivalent to saturation of $T$.
Every compatible extension containing $e\to G$ contains $W$, while
$W$ itself is a transfer system containing $V$, hence is compatible.
This proves (3).

If there is a subgroup $e<K<G$, then
$V\subsetneq W\subsetneq\mathcal C$: the first strict inclusion
is witnessed by $e\to G$, and the second by the absent arrow $K\to G$.
Part (2) makes $W$ compatible with $T$, proving that $T$ is not LSP.
If there is no such subgroup, the subgroup lattice is $\{e,G\}$,
$T=V=\id$, and its only transfer-system extensions are $V$ and
$\mathcal C$, both compatible. Thus $T$ is LSP.
Finally, a nontrivial finite group has no nontrivial proper subgroup
exactly when its order is prime. Indeed, for any $x\ne1$, either
$\langle x\rangle$ is already a nontrivial proper subgroup, or $G$ is
cyclic; a cyclic group of composite order has a proper subgroup of prime
order. This proves (4).
\end{proof}

Theorem~\ref{thm:classification} is a replacement for both conclusions
of \eqref{eq:conjecture}, rather than a change to the meaning of generation.
For nonabelian $G$, non-LSP follows under the stated hypotheses because
$|G|$ is composite, but ordinary generation is compatible only for the
specific saturated system $V$.

\begin{corollary}\label{cor:ideals}
Under the hypotheses of Theorem~\ref{thm:classification}, the compatible
greater systems are exactly
\[
 U_I=V\cup\{(K,G):K\in I\},
\]
where $I$ runs through the conjugation-invariant order ideals of the
poset of proper subgroups of $G$.
Here an order ideal is a set closed under passage to smaller subgroups;
the empty set is allowed.
This correspondence is an isomorphism of lattices, with intersection and
union as meet and join. Its only saturated members are $V$ and $\mathcal C$.
For $G=C_{p^n}$, $n\geq1$, there are exactly $n+1$ compatible greater systems.
This cyclic special case also follows from \cite[Proposition 3.2]{HML}.
\end{corollary}
\begin{proof}
Theorem~\ref{thm:classification}(2) leaves only arrows with target $G$
undetermined. If $K\to G$ belongs to a partner and $L\leq K<G$, then
$L\to K$ belongs to $V$, so composition gives $L\to G$.
Conjugation makes its set $I$ invariant.
Conversely, for such an ideal $I$, restriction of $K\to G$ to a proper
subgroup has proper target and lies in $V$. Composition before an arrow
with target $G$ uses downward closure of $I$; all other compositions
are already in $V$. Conjugation uses invariance of $I$.
Thus $U_I$ is a transfer system containing $V$, hence is compatible.
The construction is bijective and preserves intersections and unions.
If $I$ is nonempty, it contains $e$; saturation of $e\to G$ forces
every $K\to G$, so the only saturated possibilities are $I=\varnothing$
and the full proper-subgroup poset. Both are saturated.
For $C_{p^n}$, the proper subgroups form a chain with $n$ elements,
which has exactly $n+1$ order ideals.
\end{proof}

\begin{corollary}\label{cor:extremes}
For any transfer system on a nontrivial finite group with isolated $G$,
irrespective of its number of components, LSP holds if and only if $|G|$
is prime. Among transfer systems with isolated $G$, the identity system
is the unique cosaturated system, and it is bisaturated.
At the other extreme, every connected transfer system is LSP, as in
\cite[Corollary 2.15 and Theorem 3.11]{M}.
The isolated-top assertion also follows by applying
\cite[Proposition 3.17]{D} to $W$.
\end{corollary}
\begin{proof}
If $G$ is isolated, $T\subseteq V$, so $\Hull(T)\subseteq V$.
The system $W$ is compatible with $T$: for proper $A$, the conclusion
of \eqref{eq:compatibility} lies in $V$; for $A=G$, isolation forces
$B=G$, and the other antecedent is the conclusion itself.
If $|G|$ is composite, then
$\Hull(T)\subseteq V\subsetneq W\subsetneq\mathcal C$,
so $T$ is not LSP. If $|G|$ is prime, isolation forces $T=\id$,
and its only partners are $\id$ and $\mathcal C$.

By \cite[Definition 2.18]{K}, cosaturated means generated by arrows
with target $G$. An isolated top permits only $G\to G$ among these
generators, whose generated transfer system is $\id$.
The identity system is also saturated, proving the bisaturated assertion.
Finally, connectedness implies $e\to G$ by the path argument in
Theorem~\ref{thm:classification}. Its saturation is $\mathcal C$, and
compatibility forces every $K\to G$ by taking $B=e$, $C=K$ in
\eqref{eq:compatibility}; restriction then gives $\mathcal C$ as the
only partner. This also proves the cited connected-case statement.
\end{proof}

\section{Explicit counterexamples}
\begin{corollary}\label{cor:star}
For $T=P$, the ordinary generated extension is
$U=\id\cup\{(e,H):H\leq G\}$.
It is incompatible with $P$ exactly when there are subgroups
$e<C<A<G$.
Ordinary generation is compatible for every $T$ satisfying the component
hypotheses if and only if there is no such chain.
\end{corollary}
\begin{proof}
The generated-system formula follows from Theorem~\ref{thm:classification}(1).
For this star it is also the $K=e$, $H=G$ case of
\cite[Proposition A.6]{R}.
By part (3), compatibility is equivalent to $P=V$.
The arrows in $V\setminus P$ are exactly $C\to A$ with
$e<C<A<G$. This proves the equivalence.
If no such chain exists, then $P=V$, so every admissible $T$ equals $V$
and Theorem~\ref{thm:classification}(3) gives compatibility.
If such a chain exists, $T=P$ gives the required failure.
\end{proof}

Use the quaternion group
$Q_8=\langle i,j\mid i^4=1,\ i^2=j^2,\ j^{-1}ij=i^{-1}\rangle$
with $k=ij$, so its elements are $\{\pm1,\pm i,\pm j,\pm k\}$.
Put
$Z=\{\pm1\}$ and $A=\langle i\rangle=\{\pm1,\pm i\}$.
Then $e<Z<A<Q_8$. For $T=P$, the components are all five proper
subgroups and the singleton $Q_8$. Compatibility with its ordinary
generated extension would require
\[
 (e\to A)\in T,\quad (e\cap Z\to e)=(e\to e)\in U
 \quad\Longrightarrow\quad (Z\to A)\in U.
\]
The conclusion fails because $Z\ne A$ and $Z\ne e$.
This is a nonabelian counterexample to the compatibility clause.
Here the non-LSP clause is true: the compatible system $W$ lies strictly
between the hull and the complete system.

For the finite-group quantifier as printed in the source, the other clause
also has a boundary counterexample. On $C_p$ for any prime $p$, $P$ is the
identity relation on $\{e,C_p\}$ and satisfies the component hypotheses.
It is LSP, since there are exactly two transfer systems. Its ordinary
generated extension is complete and compatible.
The cyclic group $C_8$ gives another compatibility counterexample via
$e<C_2<C_4<C_8$, whereas $C_4$ illustrates that an isolated top alone does
not force failure of ordinary-generation compatibility.

\section{Coverage and verification}
The source's two assertions in Conjecture 4.2 have different obstructions.
The scope of the results is as follows.
\begin{center}\small
\begin{tabular}{>{\raggedright\arraybackslash}p{0.35\linewidth}>{\raggedright\arraybackslash}p{0.57\linewidth}}
\toprule
Source item & Coverage\\
\midrule
Conjecture 4.2, compatibility clause & Refuted-and-corrected: $Q_8$;
Theorem~\ref{thm:classification}(1)--(3). Compatibility criterion reproved;
prior proof: Mazur et al.\ (2024), Propositions 2.14, 3.5, and
DeMark et al.\ (2026), Propositions 2.21, 3.17.\\
Conjecture 4.2, non-LSP clause & Refuted-and-corrected for its printed
finite-group scope: prime-order exception; Theorem~\ref{thm:classification}(4).
Composite-order case proved; cyclic cases already in
Mazur et al.\ (2024), Theorems 4.7, 4.16 and Remark 4.18.\\
Section 4, comment immediately before Conjecture 4.2 & Not posed as a
separate assertion; isolated-top extension proved in
Corollary~\ref{cor:extremes}, without the two-component condition;
also a consequence of DeMark et al.\ (2026), Proposition 3.17.\\
Section 4, opening: connected implies LSP & Reproved; prior proof
Mazur et al.\ (2024), Corollary 2.15 and Theorem 3.11.\\
Definition 2.18 and Remark 2.19, cosaturation & Not posed; isolated-top
cosaturated and bisaturated systems classified in
Corollary~\ref{cor:extremes}.\\
Conjecture 4.1, bisaturated shortest paths & Not covered.\\
Section 3.3, final paragraph: Frobenius-width problem & Answered: yes;
prior proof cited:
Chakravarthy--Ghosh (2026), Theorem 3.2; not a consequence of our argument.\\
Remark 2.13, nonabelian complexity program & Prior proof cited for dihedral
groups: Chakravarthy--Ghosh (2026), Theorem 4.11; other families not covered.\\
All compatible partners and universal star criterion & Not posed;
proved in Corollaries~\ref{cor:ideals} and \ref{cor:star}.\\
\bottomrule
\end{tabular}
\end{center}

For the Frobenius groups $F_{p^n}=\operatorname{AGL}(1,p^n)$,
\cite[Theorem 3.2]{CG} gives
$\mathrm w(F_{p^n})=\Omega(p^n-1)+\tau(n)$, where $\Omega$ counts
prime factors with multiplicity and $\tau$ counts positive divisors.
In particular, $\mathrm w(F_p)=1+\Omega(p-1)$ answers the prime-field
part of the problem posed in \cite[Section 3.3, p.~13]{K}.
For an odd prime $p$, \cite[Theorem 4.11]{CG} gives
$\mathrm c(D_{p^n})=\lfloor3n/2\rfloor+1$, where $D_{p^n}$ has
order $2p^n$. These are cited results, rather than additional claims of
the transfer-system argument above.

The pinned
\href{https://github.com/the-omega-institute/trureturing/blob/8acba483bc80695fef00c6b8c491e2ae16d5145f/D5/S3/FiniteGroups/KlandermanConjecture42Refutation.lean}{Lean source file}
contains the declaration \texttt{result} in the namespace
\begin{center}\small
\texttt{D5.S3.FiniteGroups.KlandermanConjecture42Refutation}.
\end{center}
Its statement is the negation of the full universal conjunction in
\eqref{eq:conjecture}; its witness is the $Q_8$ system above.
The same file establishes the strict quaternion subgroup chain,
the two actual graph components, top isolation, and the exact ordinary
generated relation. The general classification in
Theorem~\ref{thm:classification}, including its LSP equivalence, has the
written proof in Section 2; it is distinct from that Lean declaration's scope.

The accompanying exact finite checker verifies the star examples, the
component classification, all isolated-top and connected systems, and the
order-ideal correspondence on $Q_8$, $C_2$, $C_4$, $C_8$, $C_6$,
$C_2\times C_2$ and $S_3$.
The proof for arbitrary finite $G$ is the written argument of
Theorem~\ref{thm:classification} and its corollaries. The source package includes the checker,
its mathematical conventions and its complete finite output.

\begin{thebibliography}{9}
\bibitem{K}
S. Klanderman, C. Lewis, H. Monson, K. Shibata and D. Van Niel,
\emph{Characterizing Transfer Systems for Non-Abelian Groups},
\href{https://arxiv.org/abs/2511.13439v2}{arXiv:2511.13439v2},
23 April 2026; first submitted 17 November 2025.

\bibitem{HML}
M. A. Hill, J. Meng and N. Li,
\emph{Counting compatible indexing systems for $C_{p^n}$},
\href{https://arxiv.org/abs/2212.11250v1}{arXiv:2212.11250v1},
21 December 2022. Proposition 3.2 uses the order (additive, multiplicative),
the reverse of the ordered pair in this note.

\bibitem{M}
K. Mazur, A. M. Osorno, C. Roitzheim, R. Santhanam, D. Van Niel and
V. Zapata Castro,
\emph{Uniquely compatible transfer systems for cyclic groups of order
$p^r q^s$}, Topology and its Applications \textbf{376} (2025), 109443.
\href{https://doi.org/10.1016/j.topol.2025.109443}{doi:10.1016/j.topol.2025.109443}.
The version used for Lemma 3.6 is
\href{https://arxiv.org/abs/2401.13523v2}{arXiv:2401.13523v2}, 21 May 2024.
The connected-case and cyclic two-component results also appear as
Corollary 2.15, Theorem 3.11, Lemma 4.14 and Theorem 4.16 in
\href{https://arxiv.org/abs/2401.13523v1}{v1}, 24 January 2024.

\bibitem{CG}
B. Chakravarthy and S. Ghosh,
\emph{Minimal generating sets of transfer systems for more non-Abelian groups},
\href{https://arxiv.org/abs/2605.04780v1}{arXiv:2605.04780v1}, 6 May 2026.

\bibitem{D}
D. DeMark, M. A. Hill, Y. Kamel, N. Niu, K. Stoeckl, D. Van Niel and G. Yan,
\emph{Maximal compatibility of disklike $G$-transfer systems},
\href{https://arxiv.org/abs/2604.00335v1}{arXiv:2604.00335v1}, 1 April 2026.

\bibitem{R}
J. Rubin,
\emph{Characterizations of equivariant little discs and linear isometries operads},
\href{https://arxiv.org/abs/1903.08723v1}{arXiv:1903.08723v1}, 20 March 2019,
Proposition A.6, p.~38.
The corresponding statement is Corollary B.8 in
\href{https://arxiv.org/abs/1903.08723v2}{v2}, 10 July 2019,
retitled \emph{Characterizations of equivariant Steiner and linear
isometries operads}.
Published as \emph{Detecting Steiner and linear isometries operads},
Glasgow Mathematical Journal \textbf{63} (2021), 307--342,
\href{https://doi.org/10.1017/S001708952000021X}{doi:10.1017/S001708952000021X}.
\end{thebibliography}

\bigskip\noindent\small
\textcopyright\ 2026 Haobo Ma and Wenlin Zhang.
Licensed under \href{https://creativecommons.org/licenses/by/4.0/}{Creative
Commons Attribution 4.0 International (CC BY 4.0)}.
\end{document}
